\documentclass[10pt]{article}
\usepackage[margin=20mm]{geometry}
\usepackage{amsmath,amssymb,mathtools,fontspec,unicode-math,xurl,hyperref}
\IfFontExistsTF{DejaVu Serif}{\setmainfont{DejaVu Serif}}{\setmainfont{Latin Modern Roman}}
\setmathfont{latinmodern-math.otf}
\setlength{\parindent}{0pt}\setlength{\parskip}{6pt}\sloppy
\title{Scientific initial review of
Hessian estimates for the sigma-2 equation
in dimension four}
\author{}
\date{26 October 2026}
\begin{document}
\maketitle
=3em\par
\section*{Identity, scope, and judgment}
This review concerns Ravi Shankar and Yu Yuan, Hessian estimates for the sigma-2 equation in dimension four, arXiv:2305.12587v2, submitted 3 December 2024, 25 PDF pages. The frozen journal identity is Annals of Mathematics 201 (2025), 489–513, DOI 10.4007/annals.2025.201.2.4. All seven manifest file bindings were independently checked. The reviewed input PDF SHA-256 is\par
00000TOKEN. The bound author TeX has 1,170 lines and SHA-256\par
00001TOKEN.\par
Judgment: zero Major findings and one Minor finding. The general, zero-order-dependent stability proof in Lemma 6.1 evaluates its operator at a function that does not touch the approximating solution. A complete counterexample to that displayed inequality is given below, together with its repair. The actual application is to a pure Hessian operator, so the omitted constant does not invalidate that application. No demonstrated failure of the main Hessian estimates was found. This is a version-limited initial review, not a certificate of whole-paper correctness.\par
All 25 pages were read; source pages 2, 5, 9, 11, 14, 16, 17, 19, 21, and 22 were visually inspected to resolve notation and location. The TeX was inspected statically; the inactive comment blocks at lines 275–281, 307–330, and 369–393 were excluded from the active proof. No author code was executed. An author-hosted journal-formatted PDF was subsequently acquired for selected comparisons, including the exact Lemma 6.1 passage; this does not replace the manifest input or constitute a full review of the later text.\par
The dispatcher reported the account Minhao Liu’s Group—Academic Researchers, Chat mode, webpage model 6Pro. Those are dispatcher observations. The visible assistant performed the scientific reading, proofs, computations, static inspection, compilation, and page inspection.\par
\section*{F01: the contact value in the general stability argument}
Equation reference: eq:wrong. Location and claim. Lemma 6.1, PDF page 22, author-source lines 1012–1032, especially line 1030. The lemma concerns continuous elliptic equations \(F(D^2u,Du,u,x)=0\) and uniformly convergent viscosity solutions \(u_k\). In its proof, after setting \[Q_\varepsilon=Q+\varepsilon|x-x_0|^2-\varepsilon^4,\] the maximum of \(u_k-Q_\varepsilon\) is positive and occurs at \(x_\varepsilon\). The proof then asserts \[
 0\le F(D^2Q_\varepsilon(x_\varepsilon),DQ_\varepsilon(x_\varepsilon),
                 Q_\varepsilon(x_\varepsilon),x_\varepsilon).\] The test function touching \(u_k\) is instead \(Q_\varepsilon+m_\varepsilon\), where \(m_\varepsilon=(u_k-Q_\varepsilon)(x_\varepsilon)>0\). Its value is \(u_k(x_\varepsilon)\), not \(Q_\varepsilon(x_\varepsilon)\).\par
Complete check under the stated assumptions. In dimension \(n=4\), on \(B_1(0)\), take \[F(M,p,z,x)=(\operatorname{tr}M)^5+z,\qquad u_k=u=0,\qquad Q=0,\quad x_0=0.\] This \(F\) is smooth and elliptic in the monotonicity sense: for \(N\ge0\), \(F(M+N,p,z,x)\ge F(M,p,z,x)\), strictly if \(N\ne0\). It is not uniformly elliptic, which the lemma does not require. No monotonicity in \(z\) is stated there. The zero functions solve the equation classically and in the viscosity sense: an upper test has nonnegative Hessian at contact, and a lower test has nonpositive Hessian. Convergence is uniform and \(Q\ge u\) with equality.\par
Here \(x_\varepsilon=0\) and \(m_\varepsilon=\varepsilon^4\). For \(\varepsilon<1/4\) the annulus estimate in the proof holds. Taking \(\varepsilon=1/65536=1/[2(2n)^5]\) gives \[F(2\varepsilon I,0,-\varepsilon^4,0)
       =(8\varepsilon)^5-\varepsilon^4
       =-\frac{\varepsilon^4}{2}<0,\] contradicting [eq:wrong] (\#eq:wrong). In fact it is negative for every \(0<\varepsilon<(2n)^{-5}\), so choosing a still smaller perturbation does not fix the step. The genuine touching test is \(\varepsilon|x|^2\); its value gives \((8\varepsilon)^5>0\), as the viscosity condition requires. Nine exact rational checks for \(2\le n\le10\) reproduce this distinction; the dimension-four wrong value is \(-1/36893488147419103232\).\par
Repair and counterevidence. Use \(\phi_\varepsilon=Q_\varepsilon+m_\varepsilon\) and write \[0\le F(D^2Q_\varepsilon(x_\varepsilon),DQ_\varepsilon(x_\varepsilon),
                    u_k(x_\varepsilon),x_\varepsilon).\] Since \(Q\ge u\), \(0<m_\varepsilon\le\left\lVert u_k-u\right\rVert_{L^\infty}+\varepsilon^4\to0\) after the indicated choice of \(k(\varepsilon)\). Also \(x_\varepsilon\to x_0\), and continuity now gives the required inequality for \(u\). The lower-test argument is analogous. Thus this is not a counterexample to the stability lemma.\par
The operator actually needed in the compactness argument has no \(z\) or gradient dependence. Adding \(m_\varepsilon\) therefore leaves that operator unchanged. Consequently no broken dependency from this error to Theorem 1.1, the higher-dimensional conditional estimate, or the four-dimensional viscosity regularity result has been demonstrated. The finding is Minor, concerning the overgeneral form of the written proof. If an additional properness condition in the zero-order variable were intended, it should be stated; the repair above needs no such extra assumption. The same unshifted zero-order value was directly checked in the author-hosted journal-formatted copy, printed page 510 (PDF page 22).\par
\section*{Jacobi inequality: definitions, algebra, and limiting cases}
Write \(H=D^2u\), \(s=\operatorname{tr}H>0\), and \(f_i=s-\lambda_i\). From \(\sigma_2(H)=1\), \[s^2=2+|H|^2,\qquad f_i>0,\qquad
 \sum_i f_i^2=(n-1)s^2-2=:\mathcal F.\] Thus the linearization is elliptic throughout the positive branch. The negative branch is handled by \(u\mapsto-u\), rather than silently included in the positive-branch argument. For \(n\ge5\) the stated dynamic hypothesis is \(\lambda_{\min}/s\ge-c_n\), where \[c_n=\frac{\sqrt{3n^2+1}-n+1}{2n}.\] For \(n=4\), the admissible-cone inequality \(\lambda_{\min}/s>-1/2\) supplies the corresponding condition automatically.\par
I redifferentiated \(\sigma_2=1\) and regrouped all symmetric third derivatives. Put \(b=\log s\), \(\delta=1+\varepsilon\), and \(t_i=(u_{11i},\ldots,u_{nni})\). Then \(Df\cdot t_i=0\) and \[s\bigl(\Delta_F b-\varepsilon|\nabla_F b|^2\bigr)
 =6\sum_{i<j<k}u_{ijk}^2+\sum_i
 \left[3|t_i|^2-2t_{i,i}^2-
       \left(1+\delta\frac{f_i}{s}\right)
                \left(\sum_j t_{i,j}\right)^2\right].\] This accounts for repeated and three-distinct-index terms; ignoring their multiplicities would change the result.\par
Let \(E\) and \(L\) be the orthogonal projections of \(e_i\) and \((1,\ldots,1)\) onto \(Df^\perp\). Direct projection gives \[|E|^2=1-\frac{f_i^2}{\mathcal F},\quad
 |L|^2=1-\frac{2(n-1)}{\mathcal F},\quad
 E\cdot L=1-\frac{(n-1)sf_i}{\mathcal F}.\] The nontrivial restriction is \(3I-2E\otimes E-\eta L\otimes L\) with \(\eta=1+\delta f_i/s\). Its trace is \(6-2|E|^2-\eta|L|^2\), and its determinant is \[d_i=9-6|E|^2-3\eta|L|^2
           +2\eta\bigl(|E|^2|L|^2-(E\cdot L)^2\bigr).\] An exact simplification, retaining favorable terms, is \[\begin{aligned}
 \frac{\mathcal F}{sf_i}d_i
 &=q_\delta(y)+\frac{6\delta}{s^2}
                   +\frac{2(n-1)\eta}{sf_i},\qquad y=\frac{f_i}{s},\\
 q_\delta(y)
 &=(n-1)(4-3\delta)
    +(6-2n+4(n-1)\delta)y-2n\delta y^2.\end{aligned}\] For \(\delta=1\), its roots are \(y_\pm=(n+1\pm\sqrt{3n^2+1})/(2n)\), and \(q_1(y)=2n(y_+-y)(y-y_-)\). The coefficient of \(\varepsilon\) is \(-3(n-1)+4(n-1)y-2ny^2>-3(n-1)\) for \(0<y<2(n-1)/n\). Therefore the paper’s choice \[\varepsilon=
 \frac{-2ny_-}{3(n-1)}
       \left(y_+-\frac{f_n}{s}\right)\ge0\] makes \(q_\delta\ge0\) for every \(i\). The trace bound and \(\delta\le3n/[2(n-1)]\) give positivity of the other invariant. Dependent \(E,L\) are covered by continuity; at a round spectrum the remaining eigenvalues can also be read directly. The case \(\varepsilon=0\) at the higher-dimensional dynamic boundary is allowed, not discarded by an implicit strict inequality.\par
The four-dimensional degenerating family \[(\lambda_1,\lambda_2,\lambda_3,\lambda_4)
       =(a,a,a,-a+1/(3a))\] has \(\sigma_2=1\) and \(\varepsilon=1/(18a^2+3)\to0\). Nevertheless \(f_4\lambda_4^2/s^3\to3/8\). This directly checks why the gradient term in the next argument rescues the regime where a fixed positive Jacobi constant would fail. For dimension three the separate \(\varepsilon=1/3\) calculation uses the sharp negative eigenvalue bound. The strict positive-edge inequality printed in Lemma 2.1 has equality at the round spectrum; its needed weak form suffices. This numerical notation issue is not counted as a scientific finding.\par
\section*{Doubling, scaling, and the dependence of constants}
The maximum-principle test on pages 10–13 combines \(2\log\rho\), \(\alpha(x\cdot Du-u)\), \(\beta|Du|^2/2\), and the logarithm of the truncated excess of \(b=\log\Delta u\) over its \(B_1\) maximum. The boundary term tends to \(-\infty\). If the excess is below the truncation threshold, the desired bound is immediate; otherwise differentiation is legitimate at the interior maximum.\par
For \(\lambda_{\min}/s\le-c_n/2\), the term \(f_n\lambda_n^2\ge(c_n/2)^2s^3\) controls the bad gradient terms after the first-derivative equation is substituted. For \(\lambda_{\min}/s\ge-c_n/2\), the Jacobi coefficient has a fixed positive lower bound. Splitting according to whether the coordinate in the largest-eigenvalue direction has square greater than \(1/n\) yields the two alternatives used in the paper: either the \(\alpha\) term is effective, or a transverse direction has \(f_k\ge(1-1/\sqrt2)s\) and a definite cutoff derivative. The possible equality in the dynamic hypothesis lies in the first regime and is therefore covered.\par
With \(\Gamma\) the stated gradient-size parameter, the choices \[\gamma=\frac1{2C},\qquad
 \alpha=\frac{\gamma^4}{\Gamma},\qquad
 \beta=\frac{\gamma^6}{\Gamma^2}\] are compatible for a sufficiently large dimensional \(C\). Indeed \(\alpha^2/\beta=\Gamma\beta/\alpha=\gamma^2\), and the final positive coefficient is proportional to \(\gamma^6(1-C\gamma-C\gamma^5)>0\). I checked these inequalities directly, including 162 exact parameter cases. They give the required exponential control of the ratio of Laplacians.\par
The translated-ball step is supported by the inclusions \(B_1(0)\subset B_{4/3}(y)\subset B_{5/3}(y)\subset B_2(0)\) for \(|y|<1/3\). The subsequent iteration uses \[u_k(x)=t^{2k}u\bigl(y+t^{-k}(x-y)\bigr),\qquad t=5/4.\] Its argument stays in \(B_3\) for \(x\in B_3\), and \(\left\lVert u_k\right\rVert_{C^1(B_3)}\le t^{2k}\left\lVert u\right\rVert_{C^1(B_3)}\). Hence a bound involving the square of that norm scales as \(t^{4k}\left\lVert u\right\rVert_{C^1}^2\). The page-14 line for \(C_k\) (source line 764) should be read with this power, rather than its printed \(t^{2k}\left\lVert u\right\rVert_{C^1}\). For example, \(u=M+|x|^2/(2\sqrt6)\) shows why a squared norm cannot be bounded uniformly by the printed linear norm. Since only finitely many steps depending on the target radius are used, the corrected estimate still has exactly the asserted dependence \(C(n,r,\left\lVert u\right\rVert_{C^1})\). This is a local quantitative correction, not a failed qualitative doubling theorem.\par
The page-11 numerical comment concerning the three-dimensional split has its comparison direction reversed. For instance \((3,3,-4/3)\) is an admissible spectrum in the first regime. The proof already handles both regimes, and no step needs that regime to be empty. This is recorded as counterevidence against escalating a numerical slip into a theorem failure.\par
\section*{Alexandrov differentiability and the compactness argument}
\subsection*{Dependency hypotheses and a complete density argument}
The proof uses two different inputs that must not be conflated. Chaudhuri–Trudinger, Theorem 2.4 [CT], gives signed Radon measures for the distributional Hessian for every \(2\le k\le n\); the stronger assumption \(k>n/2\) belongs to other conclusions of that reference. Thus its use for \(k=2,n=4\) is legitimate. The local measure argument follows by testing the dual cone with \(\mu I\) and \(\mu I-e\otimes e\), then polarizing.\par
The weighted gradient estimate in Trudinger–Wang, Corollary 3.4 [TW], applies to the constant-Hessian-measure equation as used here; it does not impose an extra negative-sign condition on \(u\). Its norm in (2.11), Theorem 3.2’s smooth Dirichlet problem, and the comparison/approximation interface were inspected. A ball is uniformly \((k-1)\)-convex; smooth boundary approximations and comparison give the stated uniform limit and interior gradient bound. This checks the hypotheses actually needed, without asserting a new proof of the entire Dirichlet theory.\par
At a point of differentiability and finite density of the Hessian measure, the mollified Taylor calculation in Lemma 4.1 is dominated by \(\left\lVert D^2u\right\rVert(B_{rt+\epsilon})/(rt+\epsilon)^n\le C\) for small radii. Its singular density tends to zero and its absolutely continuous part has a Lebesgue value. The resulting integral in \(t\) is bounded and yields the averaged second-order remainder \(o(r^2)\).\par
Two local presentations benefit from an explicit completion. At page 17, source line 842, the Taylor remainder centered at \(x\) should use \(Du(x)\) and \(D^2u(x)\), not the displayed values at 0 unless \(x\) has first been translated to 0. For the admissible quadratic \(u(y)=|y|^2/(2\sqrt6)\) and \(x=e_1/4\), the uncorrected remainder is \((y_1-1/4)/(4\sqrt6)\), whose mean absolute value on \(B_r(x)\) is \(4r/(15\pi\sqrt6)\); the correctly centered remainder is zero. The subsequent proof uses the centered version.\par
In Lemma 4.2, pages 21–22, a bad-set bound \(\varepsilon|B_r|\) by itself does not guarantee a good point within distance \(\varepsilon r\) in dimension \(n>1\). The established little-oh estimate supplies the stronger choice needed. Write \(\operatorname{Av}_B\) for the integral mean over \(B\). Fix \(0<\tau<1/2\) and take \(r\) sufficiently small that \[\operatorname{Av}_{B_r}|h|\le\frac12\tau^{n+1}r^2.\] Then \(\{|h|\ge\tau r^2\}\) has volume at most \(\tfrac12\tau^n|B_r|<|B_{\tau r}|\). Every \(y\in B_{r/2}\) therefore has a good point \(z\in B_{\tau r}(y)\). The gradient estimate (4.7), together with \(\operatorname{Av}_{B_{2r}}|h|=o(r^2)\), gives \(\operatorname{Lip}(h)\le Cr\). Consequently \(|h(y)|\le\tau r^2+C\tau r^2\). First choosing \(\tau\) small and then \(r\) small proves the desired \(\sup_{B_{r/2}}|h|=o(r^2)\), with the same hypotheses. These explicit basepoint and parameter corrections complete the argument; they are not separate formal findings or counterexamples to Alexandrov differentiability.\par
\subsection*{Savin’s local ellipticity requirements and order of choices}
At an Alexandrov point with Taylor Hessian \(H\), the upper and lower tests \(Q\pm\epsilon|x-y|^2\) imply \(\sigma_2(H)=1\) and the positive branch relation. The shifted operator used for the perturbation is \[G_H(M)=\operatorname{tr}(M+H)-\sqrt{2+|M+H|^2}.\] It is globally elliptic, since \[DG_H(M)=I-\frac{M+H}{\sqrt{2+|M+H|^2}}>0,\] and on \(|M|\le1\) it has the uniform lower bound \(1-(|H|+1)/\sqrt{2+(|H|+1)^2}>0\); its second derivatives are bounded there and \(G_H(0)=0\). Savin’s hypotheses H1, H2, H4 and Theorem 1.3 [Savin] require precisely global monotonicity and local uniform ellipticity near zero, not a global uniform ellipticity constant. This resolves a potential objection to the invocation.\par
Choose the spatial radius using the limit function’s Taylor error first, then choose the sequence index using uniform convergence. The perturbation theorem gives a neighborhood independent of that later index. The local Hessian bound combined with the moving-center doubling estimate contradicts the assumed blow-up. For \(n\ge5\), the original smooth sequence retains the dynamic condition; the approximation used solely to establish Alexandrov differentiability of the limit need not retain it. The authors do not assert unrestricted viscosity regularity in those higher dimensions.\par
For \(n=4\), the Hessian bound yields uniform ellipticity on compact subsets, allowing the indicated concavity and regularity theory to apply. The entire-solution quadratic-growth consequence uses \(u_R(x)=R^{-2}u(Rx)\); uniform local \(C^{2,\alpha}\) bounds then force oscillation of the Hessian on fixed balls to vanish as \(R\to\infty\). The general \(k\)-Hessian discussion establishes the stated differentiability/partial-regularity route, not global regularity of every admissible solution.\par
\section*{Conclusion and remaining decisions}
Item | Judgment | Repair, adoption, or remaining action\par
F01, Lemma 6.1 | Minor; displayed general inequality refuted | Add the contact constant and evaluate the zero-order variable at \(u_k(x_\varepsilon)\). Complete repair given; author adoption not requested.\par
Main estimates | No demonstrated core defect found | Jacobi algebra, dynamic boundary, two doubling regimes, Alexandrov route and local perturbation hypotheses checked as described. Not a whole-paper correctness certification.\par
Local corrections | Not counted as separate findings | Correct the norm power, Taylor basepoint and small-volume parameter choice. Explicit completions preserve the stated qualitative conclusions.\par
Version scope | Formal review: arXiv v2 | F01 also directly observed in the author-hosted typeset comparison, p. 510. Full later-text equivalence and an entire journal-version review remain unestablished.\par
Scientific decision | Initial review complete | Any editorial disposition or author response remains pending. No contact, acceptance/rejection recommendation, or publication action was taken.\par
\section*{Supporting checks}
The assistant’s own standard-library Python ran successfully (process 51432, actual exit code 0). It checked 150 exact rational dimension-four spectra; 600 exact projected-matrix trace/determinant identities and positivity cases; 150 exact symmetric third-derivative jets satisfying the differentiated equation; seven extreme spectra; 504 dimension-5–16 cases using 80-digit decimal arithmetic, including the dynamic boundary; nine exact contact counterexamples; 162 parameter choices; 18 scaling checks; 33 density-volume choices; and 120 local-ellipticity bounds. The finite checks support the general derivations above, rather than proving them by sampling. The independent script SHA-256 is\par
00002TOKEN.\par
The arXiv record lists v1 on 21 May 2023 and v2 on 3 December 2024, with no later version observed in this check. The publisher records receipt on 5 June 2023, revision on 10 September 2024, acceptance on 26 November 2024, and online publication on 12 March 2025 [publisher]. A bounded title/erratum/correction search identified no separate correction; this does not establish its nonexistence. Related variable-right-hand-side papers are separate work.\par
The author-hosted journal-formatted PDF [later] has 443,656 bytes and SHA-256\par
00003TOKEN. Its title, authors, DOI and pagination match the journal identity. Selected text comparisons concern printed pages 493, 502, 504, 509 and 510; pages 489 and 510 were additionally rendered and visually checked. In particular, the F01 formula remains on page 510. This evidence supports that specific persistence statement only. The controlled formal input remains the manifest-bound v2 PDF. The older arXiv history was inspected as metadata, not exhaustively rediffed or independently rereviewed.\par
The first report build stopped at an unavailable averaging macro, before producing a PDF. Its source and actual failure logs are preserved under `build/attempt-001`; the final source uses an explicitly defined averaging operator. The metadata command initially exceeded the Windows command-length limit without launching a process; direct file writes completed the operation, as recorded separately. The actual commands and results are retained in the adjacent evidence directory. The independent calculation was run as `python independent\_checks.py`. The report build uses the verified, previously authorized Tectonic 0.17.0 executable with `–untrusted`, `–keep-logs`, `–keep-intermediates`, `–synctex`, and `–makefile-rules`. Its executable SHA-256 is\par
00004TOKEN. No shell escape, author-code execution, software installation, optional archive export, or cloud upload was used.\par
Poppler `pdfinfo` and `pdftoppm` render the resulting PDF for inspection. The paper’s `delivery\_state.json` is finalized only after these checks. Primary web evidence, the rejected versioned dependency-fetch attempt and successful alternative source, source-image logs, static inspection, coverage, findings and provenance are retained. Classical measure differentiation, the full Dirichlet theory, Evans–Krylov theory and all cited external regularity theorems have not been independently reproved from first principles. No universal correctness or exhaustive erratum-search claim is made.\par
6 R. Shankar and Y. Yuan, Hessian estimates for the sigma-2 equation in dimension four, arXiv:2305.12587v2 (3 December 2024), bound PDF and author TeX. https://arxiv.org/abs/2305.12587. Annals of Mathematics, article metadata, DOI 10.4007/annals.2025.201.2.4. https://annals.math.princeton.edu/2025/201-2/p04. R. Shankar and Y. Yuan, author-hosted journal-formatted comparison PDF, selectively inspected. https://sites.math.washington.edu/\textasciitilde{}yuan/papers/2023aAnnals.pdf. O. Savin, Small perturbation solutions for elliptic equations, Comm. Partial Differential Equations 32 (2007), 557–578, H1, H2, H4 and Theorem 1.3. https://www.math.columbia.edu/\textasciitilde{}savin/Small\_pert.pdf. N. S. Trudinger and X.-J. Wang, Hessian measures II, Ann. of Math. 150 (1999), 579–604, (2.11), Theorem 3.2 and Corollary 3.4. https://maths-people.anu.edu.au/\textasciitilde{}wang/publications/1-Hessian-Measure2.pdf. N. Chaudhuri and N. S. Trudinger, An Alexandrov type theorem for \(k\)-convex functions, preprint dated 1 November 2004, Theorem 2.4, pp. 4–5. https://cvgmt.sns.it/media/doc/paper/1648/ct1.pdf.\par
\end{document}
